\section{可验证 Agent 的任务建模：环境、工具、状态}

\subsection{三元组视角：Environment / Tools / Verifier}
在讲座中，讲者将可验证 Agent 数据与训练对象拆成三个核心部分：环境（environment）、工具（tools）、验证器（verifier）。这不是概念分类，而是可执行系统中的三类实际约束。

\begin{knowledgebox}{Environment 的含义}
Environment 不只是 ``题目文本''，而是包含运行上下文的完整状态：代码仓库文件、依赖版本、测试脚本、系统提示词、可访问接口、历史交互痕迹等。Agent 的每个动作都可能改变该状态。
\end{knowledgebox}

\subsection{工具调用不是附属能力，而是主能力}
讲者强调，Agentic LLM 的一个核心能力是生成可被工具消费的 token 序列。也就是说，模型不仅要 ``会说''，还要 ``会调用''：参数正确、调用时机正确、错误恢复路径正确。工具链越复杂，训练数据越需要覆盖动作空间差异。

\begin{importantbox}{工具调用的三层正确性}
\begin{itemize}
    \item 语法正确：格式、字段、类型、JSON 结构合法。
    \item 语义正确：调用意图匹配任务目标，参数有业务意义。
    \item 时序正确：在正确阶段调用正确工具，避免无效循环。
\end{itemize}
\end{importantbox}

\subsection{状态转移与轨迹可验证性}
Agent 在环境中执行任务时，关键不只是最终答案，而是状态转移是否可追踪。例如代码修复任务中，``改了什么''、``为何改''、``测试是否覆盖'' 都应被记录并可验证。这就要求训练样本包含完整轨迹，而不只是终局结果。

\begin{warningbox}{只监督终局答案的风险}
如果数据只给 ``最终通过'' 样本，模型可能学到不可泛化的捷径：
\begin{itemize}
    \item 对中间步骤缺乏约束，导致不可解释行为；
    \item 面对分布外任务时，缺乏稳定的中间决策策略；
    \item 调试困难，错误定位成本飙升。
\end{itemize}
\end{warningbox}

\subsection{本章小结}
可验证 Agent 的建模核心是把问题写成 ``环境状态 + 工具动作 + 验证反馈'' 的闭环系统。训练质量取决于这三者是否被同时建模，而非只强化最终文本回答。

